\documentclass[../../../include/open-logic-section]{subfiles}

\begin{document}

\olfileid[fr]{sth}{cardinals}{cardsasords}
\olsection{Les cardinaux comme ordinaux}

En fait, notre théorie des cardinaux utilisera sans
vergogne notre théorie des ordinaux : nous définirons
les cardinaux comme certains ordinaux particuliers.
Plus précisément, nous proposerons ceci :

\begin{defn}\ollabel{defcardinalasordinal}
Si $A$ peut être bien ordonné, alors $\card{A}$
est le plus petit ordinal $\gamma$ tel que
$\cardeq{A}{\gamma}$. Pour tout ordinal $\gamma$,
on dit que $\gamma$ est un \emph{cardinal} si
et seulement si $\gamma = \card{\gamma}$.
\end{defn}

Nous venons de dire «~$A$ peut être bien ordonné~».
Comme presque toujours en mathématiques, cette tournure
modale n'est qu'une façon de parler d'une affirmation
existentielle : elle signifie qu'il existe une relation
qui met $A$ en bon ordre.

Mais la \olref{defcardinalasordinal} présente une difficulté.
Nous aimerions que \emph{tout} ensemble ait une taille,
c'est-à-dire que $\card{A}$ existe pour chaque~$A$.
Or la définition commence par une condition :
«~\emph{Si} $A$ peut être bien ordonné\ldots~».
S'il existe un ensemble $A$ qui ne peut pas être
bien ordonné, elle ne définit tout simplement aucun
objet~$\card{A}$.

Pour utiliser la \olref{defcardinalasordinal},
il nous faut donc la garantie que tout ensemble
peut être bien ordonné. Cette garantie n'est
malheureusement pas disponible dans~$\ZF$.
Si nous voulons employer
la \olref{defcardinalasordinal}, il faut donc
ajouter un nouvel axiome, par exemple :
\begin{axiom}[Bon ordre]
Tout ensemble peut être bien ordonné.
\end{axiom}
Nous examinerons dans le
\olref[choice][]{chap} si l'axiome du bon ordre
est acceptable. Dorénavant, nous l'utiliserons
simplement. Il devient alors assez facile de
démontrer que les cardinaux (au sens de la
\olref{defcardinalasordinal}) existent et se
comportent comme il faut :

\begin{lem}\ollabel{lem:CardinalsExist}
Pour tout ensemble $A$ :
\begin{enumerate}
	\item\ollabel{cardaexists}
	$\card{A}$ existe et est unique ;
	\item\ollabel{cardaapprox}
	$\cardeq{\card{A}}{A}$ ;
	\item\ollabel{cardaidem}
	$\card{A}$ est un cardinal, c'est-à-dire
	$\card{A} = \card{\card{A}}$.
\end{enumerate}
\end{lem}

\begin{proof}
Fixons $A$. D'après l'axiome du bon ordre,
il existe un bon ordre $\tuple{A, R}$.
Par le \olref[ordinals][ordtype]{thmOrdinalRepresentation},
$\tuple{A, R}$ est isomorphe à un unique ordinal
$\beta$. Donc $\cardeq{A}{\beta}$.
Par récurrence transfinie, il existe un plus petit
ordinal $\gamma$, unique, tel que
$\cardeq{A}{\gamma}$. Ainsi,
$\card{A} = \gamma$, ce qui établit
\olref{cardaexists} et \olref{cardaapprox}.
Pour \olref{cardaidem}, remarquons que
si $\delta \in \gamma$, alors
$\cardless{\delta}{A}$, par le choix de
$\gamma$ ; de plus,
$\cardless{\delta}{\gamma}$, puisque
l'équipotence est une relation d'équivalence
(\olref[sfr][siz][equ]{equinumerosityisequi}).
Donc $\gamma = \card{\gamma}$.
% Par l'absurde, supposons qu'il existe
% $\delta \in \gamma$ tel que
% $\cardeq{\gamma}{\delta}$. Comme
% $\cardeq{A}{\gamma}$, la transitivité de
% l'équipotence donnerait
% $\cardeq{A}{\delta}$, contrairement à la
% minimalité de $\gamma$.
\end{proof}

Le résultat suivant garantit le principe de Cantor,
et davantage encore. (Les cardinaux héritent de
l'ordre des ordinaux : $\cardfont{a} <
\cardfont{b}$ si et seulement si
$\cardfont{a} \in \cardfont{b}$. Dans son
énoncé, nous emploierons des lettres gothiques
pour les objets dont nous savons qu'ils sont
des cardinaux. C'est un usage assez répandu.
Une autre convention courante consiste à employer
des lettres grecques, puisque les cardinaux sont
des ordinaux, en les choisissant au milieu de
l'alphabet, par exemple $\kappa, \lambda$.) :
\begin{lem}\ollabel{lem:CardinalsBehaveRight}
Pour tous ensembles $A$ et $B$ :
\begin{align*}
	\cardeq{A}{B}
	&\text{ ssi } \card{A} = \card{B}\\
	\cardle{A}{B}
	&\text{ ssi } \card{A} \leq \card{B}\\
	\cardless{A}{B}
	&\text{ ssi } \card{A} < \card{B}
\end{align*}
\end{lem}

\begin{proof}
Nous démontrerons de gauche à droite la deuxième
équivalence (les autres cas sont semblables et
laissés en exercice). Considérons le diagramme :
\begin{center}
	\begin{tikzpicture}
	\node (nodea) {$A$};
	\node[right = 6em of nodea] (nodeb) {$B$};
	\node[below = 2em of nodea] (nodecarda) {$\card{A}$};
	\node[below = 2em of nodeb] (nodecardb) {$\card{B}$};
	\draw[->] (nodea)--(nodeb);
	\draw[<->] (nodea)--(nodecarda);
	\draw[<->] (nodeb)--(nodecardb);
	\draw[->, dashed] (nodecarda)--(nodecardb);
\end{tikzpicture}
\end{center}
Les flèches à deux pointes représentent des
!!{bijection}s, dont l'existence est garantie
par le \olref{lem:CardinalsExist}. Supposer
$\cardle{A}{B}$, c'est supposer qu'il existe
!!a{injection} $A\to B$. En suivant les flèches
de $\card{A}$ à $A$, puis à $B$, enfin à
$\card{B}$, nous obtenons !!a{injection}
$\card{A} \to \card{B}$ (la flèche en pointillés).
\end{proof}\noindent
Nous pouvons aussi réutiliser
le \olref{lem:CardinalsBehaveRight} pour redémontrer
le théorème de Schröder--Bernstein : si
$\cardle{A}{B}$ et $\cardle{B}{A}$,
alors $\cardeq{A}{B}$. Nous l'avons énoncé
dans le
\olref[sfr][siz][sb]{thm:schroder-bernstein},
mais démontré d'abord, non sans effort, dans la
\olref[sfr][infinite][card-sb]{sec}.
Voici maintenant :

\begin{proof}[Nouvelle démonstration de Schröder--Bernstein]
Si $\cardle{A}{B}$ et $\cardle{B}{A}$,
alors $\card{A} \leq \card{B}$ et
$\card{B} \leq \card{A}$ par le
\olref{lem:CardinalsBehaveRight}. Donc
$\card{A} = \card{B}$ et
$\cardeq{A}{B}$ par trichotomie et le
\olref{lem:CardinalsBehaveRight}.
\end{proof}
\noindent
Bien que cette démonstration soit très simple,
elle repose implicitement à la fois sur le
remplacement (pour obtenir
le \olref[ordinals][ordtype]{thmOrdinalRepresentation})
et sur l'axiome du bon ordre (pour garantir
le \olref{lem:CardinalsBehaveRight}).
La démonstration de la
\olref[sfr][infinite][card-sb]{sec}, en revanche,
était beaucoup plus autonome (elle peut même
être menée dans~$\Zminus$).

\end{document}
